\chapter{三相记忆的最小模型：h—E—Θ}
\chaptervoice{持续记忆不是“短期缓存加长期数据库”。本章把记忆区分为当前可控状态 $h$、带时序与结果的经历轨迹 $E$、能够解释和生成轨迹的结构 $\Theta$。三者不是三个孤立存储器，而是同一认知结构在不同时间尺度和操作方式下的相态。}

\conceptfigure{images/fig-memory.png}{三相记忆的六向流转：当前状态 $h$、经历轨迹 $E$ 与生成结构 $\Theta$。}

\section{为什么短期/长期二分不足}

短期/长期二分只按保持时间分类，无法区分两种都可长期保存但功能不同的对象：一次具体失败经历与从许多经历中提取的规律。前者需要保留时序、情境和结果，后者需要支持解释、预测和生成。

因此采用
\begin{equation}
 h=State,\qquad E=Trajectory,\qquad
 \Theta=Generator.
 \label{eq:ch23-triad}
\end{equation}
三词分别回答“当前正在发生什么”“过去怎样发生”“什么结构会生成这类轨迹”。

这种划分是功能与动力学划分，不要求对应三个物理数据库。一个载体可以实现多相，不同载体也可共同实现一相；关键在于状态、轨迹和生成器之间能否完成受控转换。

\section{$h$：当前活动状态}

$h_t$ 是时刻 $t$ 可被当前认知和控制直接读取、组合与更新的有限状态。可写为
\begin{equation}
 h_t=(CSO_t^{act},Goal_t^{act},Context_t,
      Prediction_t,Control_t,Residual_t).
 \label{eq:ch23-h}
\end{equation}
它不是所有隐状态或全部长期知识，只是当前工作集的结构化投影。

$h$ 具有高更新率、容量限制和强任务相关性。长期结构通过句柄或预测先验进入 $h$，而不是整体加载。超过容量会产生干扰、遗失上下文或错误绑定，因此调度与折叠是必要功能。

一次模型上下文可实现部分 $h$，但 $h$ 还包含目标、权限、工具状态和现实残差。把它等同提示词会丢失行动闭环与恢复语义。

\section{$E$：带情境的经验轨迹}

$E$ 保存状态—行动—结果随时间组织的经历片段：
\begin{equation}
 e_k=(t_k,h_k,u_k,y_{k+1},r_k,
      Context_k,Provenance_k).
 \label{eq:ch23-e}
\end{equation}
它保留“在什么条件下怎样做、发生了什么”，而不只保存结论。

轨迹边界由事件分段、目标变化或异常决定。片段太细会失去因果跨度，太粗会混合多个机制。写入还需记录观测与推断的区别，避免把内部预测伪装成已发生事实。

$E$ 不是不可变录像。新证据可补充解释和链接，但原始事件与后验解释应版本分离。否则回忆更新会改写历史证据，系统无法审计自己为何形成某项结构。

\section{$\Theta$：生成结构}

$\Theta$ 是从经历中形成、用于解释当前和生成未来预测的相对稳定结构。它可包含世界关系、技能、因果模型、概念、策略先验和表示变换，而不仅是传统 semantic memory。

形式上，$\Theta$ 参数化条件生成分布或状态转移：
\begin{equation}
 p_\Theta(h_{t+1},y_{t+1}\mid
          h_t,u_t,Context_t).
 \label{eq:ch23-theta}
\end{equation}
具体实现可以是规则、图、程序、参数模型或混合结构；生成器是功能身份，不限于概率神经网络。

$\Theta$ 更新慢于 $h$，但不是永恒真理。它必须保存适用域、置信度和来源，并在持续残差下修正。把它等同静态知识库会失去预测和技能生成能力。

\section{State—Trajectory—Generator 的统一关系}

三相构成不同数学类型：$h_t\in\mathcal H$ 是状态点，$E$ 是路径空间 $\mathcal H^{[t_0,t_1]}$ 上的片段集合，$\Theta$ 是路径分布或变换族的生成参数。点、路径与生成器不能用同一个“向量存储”概念替代。

统一动力学可概括为
\begin{equation}
 \begin{aligned}
 h_{t+1}&=F(h_t,u_t,y_{t+1};\Theta_t),\\
 E_{t+1}&=\mathcal A(E_t,h_t,u_t,y_{t+1}),\\
 \Theta_{t+1}&=\mathcal C(\Theta_t,E_{W_t}).
 \end{aligned}
 \label{eq:ch23-unified}
\end{equation}
$\mathcal A$ 是轨迹累积，$\mathcal C$ 是受门控巩固。

三相彼此依赖：没有 $\Theta$，$h$ 缺少预测先验；没有 $E$，结构更新失去可审计证据；没有 $h$，记忆不能重新进入当前行动。三相系统的质量应按闭环转换评价，而非单看某个存储检索率。

\section{$h\rightarrow E$：经历写入}

从 $h$ 到 $E$ 是把当前闭环片段编码为带时间、行动和结果的事件。写入算子
\begin{equation}
 e_t=W_{hE}(h_t,u_t,y_{t+1},
            Goal_t,Authority_t).
 \label{eq:ch23-h-to-e}
\end{equation}
只有在结果或明确中间状态出现后，事件才完整；计划不能预先写成成功经历。

并非每个瞬时状态都永久保存。新颖性、目标价值、残差、风险和未来可复用性共同控制写入。低价值连续信号可保存包络，高风险事件则应保留原始证据和不可篡改审计。

写入失败会造成生命周期断裂：系统执行过行动却没有轨迹，之后无法归因。重复写入则会把同一事件误作多次证据，因此需要事件身份和幂等键。

\section{$E\rightarrow h$：情境召回}

召回不是把过去原样播放，而是根据当前查询从轨迹中选取、重建并投影相关片段：
\begin{equation}
 \tilde h_t
 =R_{Eh}(E;Query_t,Context_t,Budget_t).
 \label{eq:ch23-e-to-h}
\end{equation}
$\tilde h_t$ 进入工作区前应携带来源时间、相似度和不确定度。

召回需防止情境错配。过去在不同身体、权限或环境下成功的轨迹不能直接作为当前行动依据。系统应比较适用条件并把差异显式进入 $h$。

召回也会改变当前解释，但不应覆盖事件原件。新理解可作为注释和链接写回 $E$，保留“当时记录”与“现在解释”的差别。

\section{$E\rightarrow\Theta$：结构巩固}

巩固从多条轨迹中提取可重复关系、技能或生成规律。更新可写为
\begin{equation}
 \Theta^{+}
 =\arg\min_{\Theta'}
 \left[
 \sum_{e\in E_W}\ell(e;\Theta')
 +\lambda\Omega(\Theta',\Theta)
 \right],
 \label{eq:ch23-e-to-theta}
\end{equation}
$\Omega$ 限制对既有结构的破坏。

进入 $\Theta$ 的门槛应考虑重复性、跨情境稳定、预测改善、反例和风险。单次鲜明经历可以提高候选优先级，却不应无条件改写一般规律。高风险反例可先建立例外边界，再等待更多证据。

巩固不是简单压缩文本。它需产出可用于生成、预测或控制的关系，并保留来源轨迹，使未来失配能够追溯和撤销。

\section{$\Theta\rightarrow E$：生成与解释轨迹}

$\Theta$ 可生成预测轨迹、反事实样本或技能演练：
\begin{equation}
 \hat e_{t:t+H}
 \sim G_\Theta(h_t,Goal_t,Context_t).
 \label{eq:ch23-theta-to-e}
\end{equation}
这些是 predicted 或 simulated trajectory，必须与真实经历分区标记。

该方向也包括用当前生成结构重新解释旧轨迹，例如识别共同因果模式。但解释层不能修改原始观测；应保存模型版本，使读者知道同一事件在何种 $\Theta$ 下被怎样解释。

若模拟轨迹未经区分地再次参与巩固，会形成模型训练自己输出的闭环，逐渐脱离现实。生成样本必须有权重上限和现实锚定。

\section{$\Theta\rightarrow h$：先验、技能与结构召回}

$\Theta$ 通过先验、概念句柄、策略和预测直接塑造当前状态：
\begin{equation}
 h_t^{prior}
 =P_{\Theta h}(\Theta;Context_t,Goal_t,
               Capacity_t).
 \label{eq:ch23-theta-to-h}
\end{equation}
投影应符合 $h$ 的容量，只加载当前相关结构。

熟练技能常由 $\Theta$ 直接产生动作策略，不必先召回具体经历；概念理解也可直接约束解释。这是生成结构与情景记忆的功能区别。

先验过强会压制新证据。$h$ 必须保留现实残差通道，在预测失败时降低结构置信并触发重新认知，而不是强迫观测符合 $\Theta$。

\section{$h\rightarrow\Theta$：快速修正与紧急门}

常规学习应经过 $h\to E\to\Theta$，但明确纠错、安全事件或权威规则变更可能需要快速修正。记为
\begin{equation}
 \Theta^{+}
 =U_{fast}(\Theta,h_t)
 \quad\text{s.t.}\quad Gate_{fast}=1.
 \label{eq:ch23-h-to-theta}
\end{equation}

快速门要求高置信证据、影响范围界定、权限、版本检查和回滚。例如发现某工具凭证已泄露，可立即更新安全约束；一次普通失败则不应直接重写一般模型。

即使快速修正，也必须同步生成 $E$ 事件，记录为何越过常规巩固路径。否则系统会出现无法解释的慢结构跳变。快速通道是例外治理，不是鼓励从单次上下文直接训练长期结构。

\section*{本章结论与进入下一章的接口}

$h$ 是有限当前状态，$E$ 是带情境和结果的经历轨迹，$\Theta$ 是解释并生成轨迹的结构。六向关系分别对应经历写入、情境召回、结构巩固、轨迹生成/解释、先验投影和受门控快速修正。三相不是三个数据库，也不等同传统短期/长期记忆。下一章将从信息明确度、变化速度与保持时间解释它们为何可类比气、液、晶三种信息相态。

\begin{readercheck}
选择一次工具失败事件，写出它在 $h$ 中的当前状态、写入 $E$ 的完整轨迹、可能更新的 $\Theta$ 关系；再说明何种证据允许走 $h\to\Theta$ 快速门，以及如何避免模拟轨迹被误当真实经历。
\end{readercheck}

\cleardoublepage
